\usepackage{amsmath}
\usepackage{amsfonts}
\usepackage{xspace}
\usepackage{listings}
\usepackage{mathpartir}
\usepackage{syntax}
\usepackage{subcaption}
\usepackage{hyphenat}
\usepackage{enumerate}
\usepackage{fancyvrb}
\usepackage{framed}
\usepackage{relsize}
\usepackage{bm}
\usepackage{tcolorbox}
\usepackage{textcomp}
\usepackage{mathpartir}
\usepackage{wrapfig}
\usepackage{mathtools} % reduce vertical spacing in align environments
\usepackage{tabularx}
\usepackage{caption}
\usepackage{array}
\usepackage[noend]{algorithm2e}
\usepackage{environ}
\usepackage{nopageno}
\usepackage{multicol}
\usepackage{lineno, blindtext}
\usepackage{etoolbox}
\usepackage{adjustbox}
\usepackage{changepage} % Required for the adjustwidth environment
\usepackage{tikz}
\usetikzlibrary{arrows, automata, positioning, decorations.pathmorphing}
\usepackage{forest}
\useforestlibrary{linguistics}

%% Detokenize string is empty check
\newcommand{\ifempty}[3]{\ifstrequal{#1}{}{#2}{#3}}


% Checkmark
\newcommand{\mathscale}[2]{\displaystyle\vcenter{\hbox{\scalebox{#1}{$#2$}}}}
\newcommand{\greencheck}{\mathscale{2}{\color{pinegreen}{\checkmark}}}

\newcommand{\M}{\ensuremath\mathcal{M}}
\newcommand{\Q}{\ensuremath\mathcal{Q}}

% Define an environment for programs and formulas
\NewEnviron{sentailsP}[7]{%
  \resizebox{#7}{!}{$
    \ifempty{#3}{}{\ensuremath{#3},~}
      \left[
      \hspace{-.8em}
      \begin{minipage}{#1\textwidth}
        \footnotesize
        \begin{algorithm}[H]
          #4
        \end{algorithm}
      \end{minipage}
      \hspace{-0.09\hsize},~#5~\mathscale{1.5}{\Vdash}_{#2}~#6
      \right]_S
    $}
  }

\NewEnviron{qentailsP}[8]{%
  \resizebox{#8}{!}{$
    \ifempty{#3}{}{\ensuremath{#3},~}
      \left[
      \hspace{-.8em}
      \begin{minipage}{#1\textwidth}
        \footnotesize
        \begin{algorithm}[H]
          #4
        \end{algorithm}
      \end{minipage}
      \hspace{-0.09\hsize},~#5,~#6~\mathscale{1.5}{\Vdash}_{#2}~#7
      \right]_{\Queue}
    $}
}

\NewEnviron{rentailsP}[8]{%
  \resizebox{#8}{!}{$
    \ifempty{#3}{}{\ensuremath{#3},~}
      \left[
      \hspace{-1em}
      \begin{minipage}{#1\textwidth}
        \footnotesize
        \begin{algorithm}[H]
          #4
        \end{algorithm}
      \end{minipage}
      \hspace*{-.025\hsize},~#5,~#6~\mathscale{1.5}{\Vdash}_{#2}~#7
      \right]_R
    $}
}

\NewEnviron{plet}[5]{
  \AlgoDontDisplayBlockMarkers\SetAlgoNoEnd
  \mathscale{#4}{
    \begin{minipage}{#5\textwidth}
      $\bm{#1}~{#2}\coloneq~$ \\
      \begin{algorithm}[H]
        #3
      \end{algorithm}
    \end{minipage}
  }
}

\newtheorem{theorem}{Theorem}

% For author comments
\ifx\noeditingmarks\undefined
  \newcommand{\pgwrapper}[3]{\begingroup \color{#1} #2: #3 \endgroup}
  \newcommand{\pgwrapperb}[1]{\textbf{#1}}
\else
   \newcommand{\pgwrapperb}[1]{}
   \newcommand{\pgwrapper}[2]{}
\fi

\newsavebox{\procbox}

%\newcommand{\sa}[1]{\pgwrapper{red}{SGA}{#1}\xspace}
%\newcommand{\lef}[1]{\pgwrapper{olive}{Lef}{#1}\xspace}
%\newcommand{\sz}[1]{\pgwrapper{blue}{SZ}{#1}\xspace}
%\newcommand{\yz}[1]{\pgwrapper{orange}{YZ}{#1}\xspace}
%newcommand{\todo}[1]{\pgwrapper{red}{TODO}{#1}\xspace}

% Embedding Rocq code
\usepackage{listings,lstlangcoq,lstautogobble}
\lstdefinestyle{customcoq}{
  columns=flexible,
  mathescape=true,
  belowcaptionskip=1\baselineskip,
  breaklines=true,
  numbers=none,
  xleftmargin=\parindent,
  language=Coq,
  morekeywords={Variant, fun, Arguments, Type, cofix},
  % morekeywords={SOCKAPI,ITREE,data_at,data_at_},
  emph={%
        ITree,CTree, Pid, Gt, Lt, Eq, proc, ur
  },
  emphstyle={\bfseries},
  showstringspaces=false,
  basicstyle=\small\ttfamily,
  keywordstyle=\bfseries,
  commentstyle=\itshape\color{red!40!black},
  identifierstyle=\itshape,
  stringstyle=\color{orange},
  escapeinside={<@}{@>}
}

%% Meta
\newcommand{\spaced}[1]{\ifempty{#1}{}{\ {#1}}}
\newcommand{\spacedr}[1]{\ifempty{#1}{}{{#1}\ }}
\newcommand{\dashed}[1]{\ifempty{#1}{\textcolor{gray}{\_}}{#1}}
\newcommand{\ticl}{\texttt{ticl}\xspace}
\newcommand{\Ticl}{\texttt{Ticl}\xspace}

% Shared
\newcommand{\typed}[2]{\ensuremath{#1~\textcolor{gray}{\in~#2}}}

%% From vellvm
\newcommand{\nt}[1]{\ensuremath{\mathit{#1}}} % nonterminals
\newcommand{\tm}[1]{\ensuremath{\mathtt{#1}}} % terminals
\newcommand{\kw}[1]{\ensuremath{\mathtt{#1}}} % "informal keywords"
\newcommand{\PROP}{\ensuremath{\mathbb{P}}}
\newcommand{\Type}{\ensuremath{\mathtt{Type}}}

%% Misc
\newcommand{\fin}[1]{\mathtt{fin}\spaced{#1}}
\newcommand{\unit}{\ensuremath{\mathtt{unit}}}
\newcommand{\ttt}{\mathtt{()}}
\definecolor{dgreen}{rgb}{0.0, 0.6, 0.0} % RGB values for a darker green
\definecolor{pinegreen}{rgb}{0.0, 0.47, 0.44}
\newcommand{\cbwd}{\textcolor{blue}{$\Leftarrow$}\xspace}
\newcommand{\ciff}{\textcolor{dgreen}{$\Leftrightarrow$}\xspace}
\newcommand{\coindeq}{\stackrel{\mathsf{coind}}{=}}
\newcommand{\ictree}[2]{\ensuremath{\mathtt{ictree}_{\ifempty{#1}{}{\textcolor{gray}{#1\ifempty{#2}{}{,~#2}}}}}}
\newcommand{\bbar}{\bm{|}~}


\newcommand{\ctree}{\todo{Do not refer to ctrees we have defined the syntax from scratch}\texttt{CTree}\xspace}
\newcommand{\ctrees}{\todo{Do not refer to ctrees we have defined the syntax from scratch}\texttt{CTrees}\xspace}
\newcommand{\CTree}{\todo{Do not refer to ctrees we have defined the syntax from scratch}\texttt{ctree}\xspace}
\newcommand{\CTrees}{\todo{Do not refer to ctrees we have defined the syntax from scratch}\texttt{ctrees}\xspace}
\newcommand{\itree}{\texttt{itree}\xspace}
\newcommand{\itrees}{\texttt{itrees}\xspace}
\newcommand{\ICTree}{\texttt{ICTree}\xspace}
\newcommand{\ICTrees}{\texttt{ICTrees}\xspace}
\newcommand{\ITree}{\texttt{ITree}\xspace}
\newcommand{\ITrees}{\texttt{ITrees}\xspace}
\newcommand{\sep}{~\mid~}
\newcommand{\mkw}[1]{\mbox{\sc{#1}}}
\newcommand{\upto}[1]{\ensuremath{\mkw{#1}^{UP}}}

%% CTree/Core.v
\newcommand{\Tau}[1]{\ensuremath{\mathtt{Tau}\spaced{#1}}}
\newcommand{\Br}[2]{\ensuremath{\ifempty{#1}{\mathtt{Br}\spaced{#2}}{\mathtt{Br}\spaced{#1} \spaced{#2}}}}
\newcommand{\VisX}[3]{\ensuremath{\mathtt{Vis} \spaced{#1} \spaced{#2} \spaced{#3}}}
\newcommand{\Vis}[2]{\ensuremath{\mathtt{Vis} \spaced{#1} \spaced{#2}}}
\newcommand{\Ret}[1]{\ensuremath{\mathtt{Ret} \spaced{#1}}}
\newcommand{\stuck}{\ensuremath{\tm{\emptyset}}\xspace}

%% CTree/Sbisim
\newcommand{\sbisim}{\ensuremath{\sim}}
\newcommand{\equ}{\ensuremath{\cong}}
\newcommand{\R}{\ensuremath{\mathcal{R}}}
\newcommand{\W}{\ensuremath{\mathcal{W}}}
\newcommand{\bindc}{\ensuremath{>\!\!>\!\!=}}
\newcommand{\bindt}[1]{\ensuremath{x \gets #1\tm{;\!;}}\xspace}
\newcommand{\bind}[2]{\ensuremath{x \gets #1\tm{;\!;}~#2 ~x}\xspace}
\newcommand{\bindv}[3]{\ensuremath{#1 \gets #2\tm{;\!;}~#3}\xspace}
\newcommand{\bindn}{\tm{bind}\xspace}
\newcommand{\wellfounded}[1]{\ensuremath{\mathtt{well\_founded}~#1}}
\newcommand{\iter}[2]{\ensuremath{\tm{iter\spaced{~#1}~\spaced{#2}}}\xspace}
\newcommand{\logE}[1]{\ensuremath{\mathtt{L}_{\textcolor{gray}{#1}}}}
\newcommand{\Log}[1]{\ensuremath{\mathtt{Log} \spaced{#1}}}
\newcommand{\llog}[1]{\ensuremath{\mathtt{log} \spaced{#1}}}
\newcommand{\option}[1]{\ensuremath{\tm{option}_{#1}}\xspace}
\newcommand{\Some}[1]{\ensuremath{\tm{Some}({#1})}\xspace}
\newcommand{\None}{\ensuremath{\tm{None}}\xspace}
\newcommand{\Map}[2]{\ensuremath{\tm{Map}_{\textcolor{gray}{\tm{#1, #2}}}}}
\newcommand{\Ctx}{\ensuremath{\mathcal{M}}}
\newcommand{\itern}{\tm{iter}\xspace}
\newcommand{\trigger}[1]{\ensuremath{\tm{trigger}~#1}\xspace}
\newcommand{\branch}{\tm{branch}\xspace}
\newcommand{\brtwo}{\tm{br2}\xspace}
\newcommand{\gfptop}[1]{\ensuremath{\tm{gfp}~#1}\xspace}
\newcommand{\gfp}{\ensuremath{\tm{gfp}~}}
\newcommand{\srd}[1]{\ensuremath{\tm{Rd}~{#1}}}
\newcommand{\swr}[2]{\ensuremath{\tm{Wr}~#1~#2}}
\newcommand{\sget}{\ensuremath{\tm{get}}}
\newcommand{\sput}[1]{\ensuremath{\tm{put}~#1}}
\newcommand{\cons}{\ensuremath{\tm{::}}}
\newcommand{\app}{\ensuremath{\tm{\ ++\ }}}

%% Denotation (general)
\newcommand{\sembra}[2]{\ensuremath{\llbracket \dashed{#2} \rrbracket_{\ifempty{#1}{}{_{\textcolor{gray}{\tm{#1}}}}}}}

% Equivalence
\newcommand{\sbisimR}{\sim_\R}
\newcommand{\sbisimn}{\ensuremath{\tm{sbisim}}\xspace}

%% CTree/Interp
\newcommand{\InstrM}[2]{\ensuremath{\mathtt{InstrM}_{\textcolor{gray}{{#1} ,{#2}}}}}
\newcommand{\instrv}[3]{\ensuremath{\mathtt{instr} \spaced{#1} \spaced{#2} \spaced{#3}}}
%% Instrument ctrees
\newcommand{\instr}[2]{\ensuremath{\mathcal{J}\llbracket #2 \rrbracket_{\ifempty{#1}{}{_{\textcolor{gray}{\tm{#1}}}}}}}
\newcommand{\stateT}[2]{\ensuremath{\mathtt{stateT}\spaced{\tm{#1}}\spaced{\tm{#2}}}}

%% StImp syntax
\newcommand{\stimp}{\ensuremath{\tm{StImp}}\xspace}
\newcommand{\IAExp}{\ensuremath{\tm{AExp}}\xspace}
\newcommand{\IBExp}{\ensuremath{\tm{BExp}}\xspace}
\newcommand{\St}[2]{\ensuremath{\tm{state}_{\textcolor{gray}{#1}\ifempty{#2}{}{,#2}}}}
\newcommand{\IVar}[1]{\ensuremath{\tm{var}}~{#1}}
\newcommand{\IVal}[1]{\ensuremath{\tm{val}}~{#1}}
\newcommand{\IAdd}[2]{\ensuremath{#1 ~ + ~ #2}}
\newcommand{\ISub}[2]{\ensuremath{#1 ~ - ~ #2}}
\newcommand{\IEq}[2]{\ensuremath{#1 ~ = ~ #2}}
\newcommand{\ILt}[2]{\ensuremath{#1 ~ < ~ #2}}
\newcommand{\ILe}[2]{\ensuremath{#1 ~ \leq ~ #2}}
\newcommand{\IIf}[3]{\ensuremath{\text{\textbf{if}}~(#1)~\text{\textbf{then}}~ #2 ~\text{\textbf{else}}~ #3}}
\newcommand{\IWhile}[2]{\ensuremath{\text{\textbf{while}}~(#1)~\{#2\}}}
\newcommand{\IAssign}[2]{\ensuremath{#1~ \leftarrow ~ #2}}
\newcommand{\ISeq}[2]{\ensuremath{#1~;~#2}}
\newcommand{\IBr}[2]{\ensuremath{#1~\oplus~#2}}
\newcommand{\ISkip}{\ensuremath{\tm{skip}}}

%% MeQ syntax
\newcommand{\meQ}[1]{\ensuremath{\tm{MeQ}\ifempty{#1}{}{_{#1}}}\xspace}
\newcommand{\QPop}{\ensuremath{\tm{pop}}\xspace}
\newcommand{\QPush}[1]{\ensuremath{\tm{push} \spaced{#1}}\xspace}
\newcommand{\QWhile}[1]{\ensuremath{\text{\textbf{while}}~(\tm{true})~\{#1\}}}
\newcommand{\QRet}[1]{\ensuremath{\tm{ret}~#1}}
\newcommand{\QBind}[2]{\ensuremath{#1 \bindc #2}}

%% MeS syntax
\newcommand{\meS}[1]{\ensuremath{\tm{MeS}\ifempty{#1}{}{_{#1}}}\xspace}
\newcommand{\SWrite}[3]{\ensuremath{\tm{write} \spaced{#1} \spaced{#2} \spaced{#3}}\xspace}
\newcommand{\SRead}[2]{\ensuremath{\tm{read} \spaced{#1} \spaced{#2}}\xspace}
\newcommand{\SIf}[3]{\ensuremath{\text{\textbf{if}} ~#1~ \text{\textbf{then}}~#2~\text{\textbf{else}}~#3}\xspace}
\newcommand{\SRet}[1]{\ensuremath{\tm{ret}~#1}}
\newcommand{\SBind}[2]{\ensuremath{#1 \bindc #2}}

%% MeR syntax
\newcommand{\meR}[1]{\ensuremath{\tm{MeR}\ifempty{#1}{}{_{#1}}}\xspace}
\newcommand{\RLoop}[2]{\ensuremath{\tm{loop}} \spaced{#1} \spaced{#2} \xspace}
\newcommand{\RBr}[2]{\ensuremath{#1~\oplus~#2}}
\newcommand{\RCall}[1]{\ensuremath{\tm{call} ~#1}\xspace}
\newcommand{\RRet}[1]{\ensuremath{\tm{ret}~#1}}
\newcommand{\RBind}[2]{\ensuremath{#1 \bindc #2}}

%% Logic/Kripke.v
\newcommand{\World}[1]{\ensuremath{\mathcal{W}\ifempty{#1}{}{_{\textcolor{gray}{#1}}}}}
\newcommand{\Pure}{\ensuremath{\mathtt{Pure}}\xspace}
\newcommand{\Val}[1]{\ensuremath{\mathtt{Val} \spaced{#1}}\xspace}
\newcommand{\Obs}[2]{\ensuremath{\mathtt{Obs} \spaced{#1} \spaced{#2}}\xspace}
\newcommand{\Finish}[3]{\ensuremath{\mathtt{Finish} \spaced{#1} \spaced{#2} \spaced{#3}}\xspace}
\newcommand{\notdone}[1]{\ensuremath{\mathtt{not\_done} \spaced{#1}}}
\newcommand{\donewith}[2]{\ensuremath{\mathtt{done\_with} \spaced{#1} \spaced{#2}}}
\newcommand{\canstep}[2]{\ensuremath{\mathtt{can\_step} \spaced{#1} \spaced{#2}}}
\newcommand{\ktrans}[4]{\ensuremath{[\dashed{#1},\ \dashed{#2} ] \mapsto [ \dashed{#3},\ \dashed{#4} ]}}

%%% Logic/Syntax.v
\newcommand{\psix}{\ensuremath{\psi_\mathtt{X}}\xspace}
\newcommand{\Px}{\ensuremath{P_\mathtt{X}}\xspace}
\newcommand{\psiy}{\ensuremath{\psi_\mathtt{Y}}\xspace}
\newcommand{\phic}{\ensuremath{\varphi}\xspace}
\newcommand{\implL}{\ensuremath{\Rightarrow_L}}
\newcommand{\implR}{\ensuremath{\Rightarrow_R}}
\newcommand{\implLR}{\ensuremath{\Rightarrow_{LR}}}
\newcommand{\iffL}{\ensuremath{\Leftrightarrow_L}}
\newcommand{\iffR}{\ensuremath{\Leftrightarrow_R}}
\newcommand{\iffLR}{\ensuremath{\Leftrightarrow_{LR}}}

%%%% SBisim/Trans.v
\newcommand{\lts}{\texttt{LTS}\xspace}
\newcommand{\lbl}{\ensuremath{\mathtt{label}}\xspace}
\newcommand{\ltau}{\ensuremath{\tau}\xspace}
\newcommand{\lobs}[2]{\ensuremath{\mathtt{obs}\spaced{#1}\spaced{#2}}\xspace}
\newcommand{\lval}[1]{\ensuremath{\mathtt{val}\spaced{#1}}\xspace}
\newcommand{\lstep}[3]{#1 \xrightarrow{#2} #3\xspace}

% base predicates
\newcommand{\now}[1]{\ensuremath{\mathtt{now}\spaced{#1}}\xspace}
\newcommand{\return}[1]{\ensuremath{\mathtt{return}\spaced{#1}}\xspace}
% derived
\newcommand{\done}[1]{\ensuremath{\mathtt{done}\spaced{#1}}\xspace}
\newcommand{\doneq}[2]{\ensuremath{\mathtt{done_=}\spaced{#1}\spaced{#2}}\xspace}
\newcommand{\val}[1]{\ensuremath{\mathtt{val}\spaced{#1}}\xspace}
\newcommand{\finish}[1]{\ensuremath{\mathtt{finish}\spaced{#1}}\xspace}
\newcommand{\pure}{\ensuremath{\mathtt{pure}}\xspace}
\newcommand{\obs}[1]{\ensuremath{\mathtt{obs}\spaced{#1}}\xspace}
% inl
\newcommand{\inl}[1]{\ensuremath{\mathtt{inl}\spaced{#1}}\xspace}
\newcommand{\inr}[1]{\ensuremath{\mathtt{inr}\spaced{#1}}\xspace}

\newcommand{\AU}[2]{\ensuremath{\spacedr{#1} \tm{AU}      \spaced{#2}}\xspace}
\newcommand{\EU}[2]{\ensuremath{\spacedr{#1} \tm{EU}      \spaced{#2}}\xspace}
\newcommand{\AN}[2]{\ensuremath{\spacedr{#1} \tm{AN}      \spaced{#2}}\xspace}
\newcommand{\EN}[2]{\ensuremath{\spacedr{#1} \tm{EN}      \spaced{#2}}\xspace}
\newcommand{\AG}[1]{\ensuremath{\tm{AG}\spaced{#1}}\xspace}
\newcommand{\EG}[1]{\ensuremath{\tm{EG}\spaced{#1}}\xspace}

%% Syntactic sugar
\newcommand{\AX}[1]{\ensuremath{\tm{AX} \spaced{#1}\xspace}}
\newcommand{\EX}[1]{\ensuremath{\tm{EX} \spaced{#1}\xspace}}
\newcommand{\AF}[1]{\ensuremath{\tm{AF} \spaced{#1}\xspace}}
\newcommand{\EF}[1]{\ensuremath{\tm{EF} \spaced{#1}\xspace}}

%% Entailment
\newcommand{\entails}[3]{\ensuremath{\langle ~\dashed{#1}, \ \dashed{#2} \vDash ~\dashed{#3}~ \rangle}}
\newcommand{\entailsL}[3]{\ensuremath{\langle ~\dashed{#1}, \ \dashed{#2} \vDash_L ~\dashed{#3}~ \rangle}}
\newcommand{\entailsR}[3]{\ensuremath{\langle ~\dashed{#1}, \ \dashed{#2} \vDash_R ~\dashed{#3}~ \rangle}}
\newcommand{\entailsRR}[2]{\ensuremath{\langle~ \dashed{#1}, \ \dashed{#2} \vDash_R ~}}
\newcommand{\entailsLR}[3]{\ensuremath{\langle~ \dashed{#1}, \ \dashed{#2} \vDash_{LR} ~\dashed{#3} ~\rangle}}
%% stimp entailment
\newcommand{\sentailsL}[4]{\ensuremath{[\dashed{#2}, \ \dashed{#3} \Vdash_L ~\dashed{#4} ]_{#1}}}
\newcommand{\sentailsR}[4]{\ensuremath{[\dashed{#2}, \ \dashed{#3} \Vdash_R ~\dashed{#4} ]_{#1}}}
\newcommand{\sentailsLR}[4]{\ensuremath{[\dashed{#2}, \ \dashed{#3} \Vdash_{LR} ~\dashed{#4} ]_{#1}}}
\newcommand{\sentailsRR}[2]{\ensuremath{[\dashed{#1}, \ \dashed{#2} \Vdash_R ~ }}

%% Queue entailment
\newcommand{\qentailsL}[4]{\ensuremath{[\dashed{#1}, \ \dashed{#2},\ \dashed{#3} \Vdash_L ~\dashed{#4} ]_{\Queue}}}
\newcommand{\qentailsR}[4]{\ensuremath{[\dashed{#1}, \ \dashed{#2},\ \dashed{#3} \Vdash_R ~\dashed{#4} ]_{\Queue}}}
\newcommand{\qentailsLR}[4]{\ensuremath{[\dashed{#1}, \ \dashed{#2},\ \dashed{#3} \Vdash_{LR} ~\dashed{#4} ]_{\Queue}}}
\newcommand{\qentailsRR}[3]{\ensuremath{[\dashed{#1}, \ \dashed{#2},\ \dashed{#3} \Vdash_R ~}}

%% Sec memory entailment
\newcommand{\mentailsL}[4]{\ensuremath{[\dashed{#1}, \ \dashed{#2},\ \dashed{#3} \Vdash_L ~\dashed{#4} ]_{\Lbl}}}
\newcommand{\mentailsR}[4]{\ensuremath{[\dashed{#1}, \ \dashed{#2},\ \dashed{#3} \Vdash_R ~\dashed{#4} ]_{\Lbl}}}
\newcommand{\mentailsLR}[4]{\ensuremath{[\dashed{#1}, \ \dashed{#2},\ \dashed{#3} \Vdash_{LR} ~\dashed{#4} ]_{\Lbl}}}
\newcommand{\mentailsRR}[3]{\ensuremath{[\dashed{#1}, \ \dashed{#2},\ \dashed{#3} \Vdash_R ~}}
\newcommand{\rentailsL}[4]{\ensuremath{[\dashed{#1}, \ \dashed{#2},\ \dashed{#3} \Vdash_L ~\dashed{#4} ]_{R}}}
\newcommand{\rentailsR}[4]{\ensuremath{[\dashed{#1}, \ \dashed{#2},\ \dashed{#3} \Vdash_R ~\dashed{#4} ]_{R}}}
\newcommand{\rentailsLR}[4]{\ensuremath{[\dashed{#1}, \ \dashed{#2},\ \dashed{#3} \Vdash_{LR} ~\dashed{#4} ]_{R}}}
\newcommand{\rentailsRR}[3]{\ensuremath{[\dashed{#1}, \ \dashed{#2},\ \dashed{#3} \Vdash_R ~}}
\newcommand{\rentailsLL}[3]{\ensuremath{[\dashed{#1}, \ \dashed{#2},\ \dashed{#3} \Vdash_L ~}}

%%% StImp
\newcommand{\Get}{\ensuremath{\tm{Get}}}
\newcommand{\Put}[1]{\ensuremath{\tm{Put~{#1}}}}
\newcommand{\svar}[1]{\ensuremath{\tm{var}~{#1}}}

%%% Examples, queue
\newcommand{\Queue}{\ensuremath{\mathcal{Q}}}
\newcommand{\qpop}{\tm{pop}}
\newcommand{\qpush}[1]{\ensuremath{\tm{push}\spaced{#1}}}


%%% Examples, secure memory
\newcommand{\MS}{\ensuremath{\mathcal{M}_\mathcal{S}}}
\newcommand{\Lbl}{\ensuremath{\mathcal{S}}}
\newcommand{\stimpS}{\ensuremath{\stimp_\mathcal{S}}\xspace}
\newcommand{\sread}[2]{\ensuremath{\tm{read}_{#1}~#2}}
\newcommand{\swrite}[3]{\ensuremath{\tm{write}_{#1}~#2~#3}}
\newcommand{\siseven}[1]{\ensuremath{\tm{is\_even}~#1}}
\newcommand{\IAExpS}{\ensuremath{\tm{AExp_\mathcal{S}}}}
\newcommand{\IBExpS}{\ensuremath{\tm{BExp_\mathcal{S}}}}

%%% Examples, election
\newcommand{\Pid}{\ensuremath{\tm{PID}_n}}
\newcommand{\Enet}{\ensuremath{E_{\tm{net}}}}
\newcommand{\hnet}{\ensuremath{h_{\tm{net}}}}
\newcommand{\Msg}{\ensuremath{\tm{Msg}_n}}
\newcommand{\Mailboxes}{\ensuremath{[\Msg]_n}}
\newcommand{\esend}[2]{\ensuremath{\tm{send} \spaced{#1} \spaced{#2}}}
\newcommand{\erecv}[1]{\ensuremath{\tm{recv} \spaced{#1}}}
\newcommand{\proc}[1]{\ensuremath{\tm{proc} \spaced{#1}}}

\usepackage{tabularray}

\NewEnviron{typsyntax}[2]{%
  \begin{flushleft}
    \begin{tabular}{l l}
      $\typed{#1}{#2}~=$
      &
        $\begin{aligned}[t] % Start aligned in math mode
          \BODY
        \end{aligned}$ % End aligned and math mode
    \end{tabular}
  \end{flushleft}
}

\NewEnviron{typcosyntax}[2]{%
  \begin{flushleft}
    \begin{tabular}{l l}
      $\typed{#1}{#2}~=$
      &
        $\begin{aligned}[t] % Start aligned in math mode
          \BODY
        \end{aligned}$ % End aligned and math mode
    \end{tabular}
  \end{flushleft}
}

\NewEnviron{typfunction}[2]{%
  \begin{flushleft}
    \begin{tabular}{l}
      \ensuremath{\mathtt{#1}~\textcolor{gray}{\in~#2}} \\[.5em]
      {\begin{minipage}{\textwidth}
          $\begin{aligned} % Start aligned in math mode
            \BODY
          \end{aligned}$
        \end{minipage}
      } % End aligned and math mode
    \end{tabular}
  \end{flushleft}
}

\newcommand{\typequation}[3]{
  \begin{typsyntax}{#1}{#2}
    #3 \\
  \end{typsyntax}
}

\newcommand{\typcoequation}[3]{
  \begin{typcosyntax}{#1}{#2}
    #3 \\
  \end{typcosyntax}
}
